
\keyheading{第1题　分类（填序号）}
\begingroup\color{answerred}平行四边形：①④⑤；梯形：②③⑥。\par\endgroup
\keyheading{第2题　底为5厘米、7厘米的平行四边形及高}
\begingroup\color{answerred}\par\noindent\begin{center}{%
\begin{tikzpicture}[x=1bp,y=-1bp,color=black,draw=black,text=black]
\begin{scope}\clip (64.00,344.00) rectangle (555.00,477.00);
\hdhead{44}{多边形的认识——信息窗3}
\hdsection{1}{145}{分一分。（填序号。）}
\begin{scope}[scale=.6666667,draw=black,text=black]\draw[hdoutline] (108.000,229.000)--(197.000,229.000)--(218.000,274.000)--(129.000,274.000)--cycle;\end{scope}
\begin{scope}[scale=.6666667,draw=black,text=black]\draw[hdoutline] (249.000,273.000)--(257.000,228.000)--(318.000,228.000)--(344.000,273.000)--cycle;\end{scope}
\begin{scope}[scale=.6666667,draw=black,text=black]\draw[hdoutline] (375.000,272.000)--(391.000,227.000)--(445.000,227.000)--(463.000,272.000)--cycle;\end{scope}
\begin{scope}[scale=.6666667,draw=black,text=black]\draw[hdoutline] (498.000,271.000)--(517.000,226.000)--(594.000,226.000)--(575.000,271.000)--cycle;\end{scope}
\begin{scope}[scale=.6666667,draw=black,text=black]\draw[hdoutline] (630.000,232.000)--(660.000,194.000)--(690.000,232.000)--(660.000,270.000)--cycle;\end{scope}
\begin{scope}[scale=.6666667,draw=black,text=black]\draw[hdoutline] (723.000,225.000)--(800.000,225.000)--(794.000,269.000)--(752.000,269.000)--cycle;\end{scope}
\begin{scope}[scale=.6666667,draw=black,text=black]\node[inner sep=0pt,font=\normalsize,text=black] at (163.000,311.000) {①};\end{scope}
\begin{scope}[scale=.6666667,draw=black,text=black]\node[inner sep=0pt,font=\normalsize,text=black] at (296.000,311.000) {②};\end{scope}
\begin{scope}[scale=.6666667,draw=black,text=black]\node[inner sep=0pt,font=\normalsize,text=black] at (419.000,311.000) {③};\end{scope}
\begin{scope}[scale=.6666667,draw=black,text=black]\node[inner sep=0pt,font=\normalsize,text=black] at (547.000,311.000) {④};\end{scope}
\begin{scope}[scale=.6666667,draw=black,text=black]\node[inner sep=0pt,font=\normalsize,text=black] at (660.000,311.000) {⑤};\end{scope}
\begin{scope}[scale=.6666667,draw=black,text=black]\node[inner sep=0pt,font=\normalsize,text=black] at (764.000,311.000) {⑥};\end{scope}
\q{144}{342}{132}{13}{平行四边形：}
\q{262}{339}{28}{13}{（}
\q{459}{339}{25}{13}{）}
\q{527}{342}{80}{13}{梯形：}
\q{583}{339}{28}{13}{（}
\q{760}{339}{25}{13}{）}
\hdsection{2}{403}{按要求画一画。}
\q{104}{447}{734}{12}{在方格纸上分别画一个底是5厘米和一个底是7厘米的平行四边形，}
\q{104}{489}{734}{12}{并各画出一条高。（每个小方格的边长为1厘米。）}
\begin{scope}[scale=.6666667,draw=black,text=black]\draw[hdgrid] (103,531) rectangle (813,700);\foreach \i in {1,...,29}{\draw[hdgrid] ({103+23.666666666666668*\i},531)--({103+23.666666666666668*\i},700);}\foreach \j in {1,...,6}{\draw[hdgrid] (103,{531+24.142857142857142*\j})--(813,{531+24.142857142857142*\j});}\end{scope}
\hdsection{3}{736}{分别画出下面平行四边形底边上的高。}
\begin{scope}[scale=.6666667,draw=black,text=black]\draw[hdoutline] (123.000,860.000)--(147.000,808.000)--(227.000,808.000)--(203.000,860.000)--cycle;\end{scope}
\begin{scope}[scale=.6666667,draw=black,text=black]\draw[hdoutline] (307.000,865.000)--(332.000,782.000)--(386.000,782.000)--(361.000,865.000)--cycle;\end{scope}
\begin{scope}[scale=.6666667,draw=black,text=black]\draw[hdoutline] (471.000,785.000)--(528.000,806.000)--(528.000,882.000)--(471.000,861.000)--cycle;\end{scope}
\begin{scope}[scale=.6666667,draw=black,text=black]\draw[hdoutline] (670.000,808.000)--(755.000,808.000)--(781.000,861.000)--(696.000,861.000)--cycle;\end{scope}
\begin{scope}[scale=.6666667,draw=black,text=black]\node[inner sep=0pt,font=\normalsize,text=black] at (118.000,830.000) {底};\end{scope}
\begin{scope}[scale=.6666667,draw=black,text=black]\node[inner sep=0pt,font=\normalsize,text=black] at (333.000,881.000) {底};\end{scope}
\begin{scope}[scale=.6666667,draw=black,text=black]\node[inner sep=0pt,font=\normalsize,text=black] at (546.000,842.000) {底};\end{scope}
\begin{scope}[scale=.6666667,draw=black,text=black]\node[inner sep=0pt,font=\normalsize,text=black] at (662.000,839.000) {底};\end{scope}
\hdsection{4}{923}{画一个以$BC$为底、$AD$为高的平行四边形。}
\begin{scope}[scale=.6666667,draw=black,text=black]\draw[hdgrid] (241,967) rectangle (673,1133);\foreach \i in {1,...,20}{\draw[hdgrid] ({241+20.571428571428573*\i},967)--({241+20.571428571428573*\i},1133);}\foreach \j in {1,...,7}{\draw[hdgrid] (241,{967+20.75*\j})--(673,{967+20.75*\j});}\end{scope}
\begin{scope}[scale=.6666667,draw=black,text=black]\draw[hdoutline] (382.000,1071.000)--(506.000,1071.000);\end{scope}
\begin{scope}[scale=.6666667,draw=black,text=black]\draw[hdoutline,dashed] (424.000,1009.000)--(424.000,1071.000);\end{scope}
\begin{scope}[scale=.6666667,draw=black,text=black]\draw[line width=.45pt] (424.000,1063.000)--(432.000,1063.000)--(432.000,1071.000);\end{scope}
\begin{scope}[scale=.6666667,draw=black,text=black]\node[inner sep=0pt,font=\normalsize,text=black] at (424.000,991.000) {$A$};\end{scope}
\begin{scope}[scale=.6666667,draw=black,text=black]\node[inner sep=0pt,font=\normalsize,text=black] at (382.000,1088.000) {$B$};\end{scope}
\begin{scope}[scale=.6666667,draw=black,text=black]\node[inner sep=0pt,font=\normalsize,text=black] at (506.000,1088.000) {$C$};\end{scope}
\begin{scope}[scale=.6666667,draw=black,text=black]\node[inner sep=0pt,font=\normalsize,text=black] at (424.000,1088.000) {$D$};\end{scope}

\begin{scope}[color=answerred,draw=answerred,text=answerred]\ans{190.00}{226.00}{106.67}{11}{①④⑤}
\ans{406.00}{226.00}{100.00}{11}{②③⑥}
\draw[answerred,line width=.8pt] (99.333,450.667) -- (177.333,450.667) -- (208.667,386.000) -- (130.667,386.000) -- cycle;
\draw[answerred,line width=.7pt,dashed] (130.667,386.000) -- (130.667,450.667);
\draw[answerred,line width=.7pt] (136.000,450.667) -- (136.000,445.333);
\draw[answerred,line width=.7pt] (136.000,445.333) -- (130.667,445.333);
\ans{134.67}{413.00}{20.00}{9}{高}
\ans{136.00}{452.67}{66.67}{9}{$\vphantom{0}5\mathrm{cm}$}
\draw[answerred,line width=.8pt] (301.333,450.000) -- (411.333,450.000) -- (443.333,385.333) -- (333.333,385.333) -- cycle;
\draw[answerred,line width=.7pt,dashed] (333.333,385.333) -- (333.333,450.000);
\draw[answerred,line width=.7pt] (338.667,450.000) -- (338.667,444.667);
\draw[answerred,line width=.7pt] (338.667,444.667) -- (333.333,444.667);
\ans{337.33}{412.33}{20.00}{9}{高}
\ans{356.00}{452.67}{66.67}{9}{$\vphantom{0}7\mathrm{cm}$}
\draw[answerred,line width=.7pt,dashed] (135.333,573.333) -- (91.366,553.041);
\draw[answerred,line width=.7pt] (89.131,557.883) -- (93.973,560.118);
\draw[answerred,line width=.7pt] (93.973,560.118) -- (96.208,555.276);
\ans{117.35}{557.85}{20.00}{9}{高}
\draw[answerred,line width=.7pt,dashed] (221.333,521.333) -- (221.333,576.667);
\draw[answerred,line width=.7pt] (226.667,576.667) -- (226.667,571.333);
\draw[answerred,line width=.7pt] (226.667,571.333) -- (221.333,571.333);
\ans{225.33}{543.67}{20.00}{9}{高}
\draw[answerred,line width=.7pt,dashed] (314.000,574.000) -- (352.000,574.000);
\draw[answerred,line width=.7pt] (352.000,568.667) -- (346.667,568.667);
\draw[answerred,line width=.7pt] (346.667,568.667) -- (346.667,574.000);
\ans{337.00}{568.67}{20.00}{9}{高}
\draw[answerred,line width=.7pt,dashed] (503.333,538.667) -- (457.659,561.073);
\draw[answerred,line width=.7pt] (455.310,556.285) -- (460.098,553.936);
\draw[answerred,line width=.7pt] (460.098,553.936) -- (462.447,558.724);
\ans{484.50}{544.54}{20.00}{9}{高}
\draw[answerred,line width=.8pt] (254.667,714.000) -- (337.333,714.000) -- (365.333,672.667) -- (282.667,672.667) -- cycle;
\draw[answerred,line width=.7pt,dashed] (282.667,672.667) -- (282.667,714.000);
\ans{286.5}{688.0}{20}{9}{高}
\ans{366.67}{656.67}{23.33}{9}{$\vphantom{0}E$}
\end{scope}
\end{scope}\pgfresetboundingbox\path[use as bounding box] (64.00,344.00) rectangle (555.00,477.00);
\end{tikzpicture}}
\end{center}\par\endgroup
\keyheading{第3题　指定底边上的高}
\begingroup\color{answerred}\par\noindent\begin{center}{%
\begin{tikzpicture}[x=1bp,y=-1bp,color=black,draw=black,text=black]
\begin{scope}\clip (55.00,519.00) rectangle (554.00,609.00);
\hdhead{44}{多边形的认识——信息窗3}
\hdsection{1}{145}{分一分。（填序号。）}
\begin{scope}[scale=.6666667,draw=black,text=black]\draw[hdoutline] (108.000,229.000)--(197.000,229.000)--(218.000,274.000)--(129.000,274.000)--cycle;\end{scope}
\begin{scope}[scale=.6666667,draw=black,text=black]\draw[hdoutline] (249.000,273.000)--(257.000,228.000)--(318.000,228.000)--(344.000,273.000)--cycle;\end{scope}
\begin{scope}[scale=.6666667,draw=black,text=black]\draw[hdoutline] (375.000,272.000)--(391.000,227.000)--(445.000,227.000)--(463.000,272.000)--cycle;\end{scope}
\begin{scope}[scale=.6666667,draw=black,text=black]\draw[hdoutline] (498.000,271.000)--(517.000,226.000)--(594.000,226.000)--(575.000,271.000)--cycle;\end{scope}
\begin{scope}[scale=.6666667,draw=black,text=black]\draw[hdoutline] (630.000,232.000)--(660.000,194.000)--(690.000,232.000)--(660.000,270.000)--cycle;\end{scope}
\begin{scope}[scale=.6666667,draw=black,text=black]\draw[hdoutline] (723.000,225.000)--(800.000,225.000)--(794.000,269.000)--(752.000,269.000)--cycle;\end{scope}
\begin{scope}[scale=.6666667,draw=black,text=black]\node[inner sep=0pt,font=\normalsize,text=black] at (163.000,311.000) {①};\end{scope}
\begin{scope}[scale=.6666667,draw=black,text=black]\node[inner sep=0pt,font=\normalsize,text=black] at (296.000,311.000) {②};\end{scope}
\begin{scope}[scale=.6666667,draw=black,text=black]\node[inner sep=0pt,font=\normalsize,text=black] at (419.000,311.000) {③};\end{scope}
\begin{scope}[scale=.6666667,draw=black,text=black]\node[inner sep=0pt,font=\normalsize,text=black] at (547.000,311.000) {④};\end{scope}
\begin{scope}[scale=.6666667,draw=black,text=black]\node[inner sep=0pt,font=\normalsize,text=black] at (660.000,311.000) {⑤};\end{scope}
\begin{scope}[scale=.6666667,draw=black,text=black]\node[inner sep=0pt,font=\normalsize,text=black] at (764.000,311.000) {⑥};\end{scope}
\q{144}{342}{132}{13}{平行四边形：}
\q{262}{339}{28}{13}{（}
\q{459}{339}{25}{13}{）}
\q{527}{342}{80}{13}{梯形：}
\q{583}{339}{28}{13}{（}
\q{760}{339}{25}{13}{）}
\hdsection{2}{403}{按要求画一画。}
\q{104}{447}{734}{12}{在方格纸上分别画一个底是5厘米和一个底是7厘米的平行四边形，}
\q{104}{489}{734}{12}{并各画出一条高。（每个小方格的边长为1厘米。）}
\begin{scope}[scale=.6666667,draw=black,text=black]\draw[hdgrid] (103,531) rectangle (813,700);\foreach \i in {1,...,29}{\draw[hdgrid] ({103+23.666666666666668*\i},531)--({103+23.666666666666668*\i},700);}\foreach \j in {1,...,6}{\draw[hdgrid] (103,{531+24.142857142857142*\j})--(813,{531+24.142857142857142*\j});}\end{scope}
\hdsection{3}{736}{分别画出下面平行四边形底边上的高。}
\begin{scope}[scale=.6666667,draw=black,text=black]\draw[hdoutline] (123.000,860.000)--(147.000,808.000)--(227.000,808.000)--(203.000,860.000)--cycle;\end{scope}
\begin{scope}[scale=.6666667,draw=black,text=black]\draw[hdoutline] (307.000,865.000)--(332.000,782.000)--(386.000,782.000)--(361.000,865.000)--cycle;\end{scope}
\begin{scope}[scale=.6666667,draw=black,text=black]\draw[hdoutline] (471.000,785.000)--(528.000,806.000)--(528.000,882.000)--(471.000,861.000)--cycle;\end{scope}
\begin{scope}[scale=.6666667,draw=black,text=black]\draw[hdoutline] (670.000,808.000)--(755.000,808.000)--(781.000,861.000)--(696.000,861.000)--cycle;\end{scope}
\begin{scope}[scale=.6666667,draw=black,text=black]\node[inner sep=0pt,font=\normalsize,text=black] at (118.000,830.000) {底};\end{scope}
\begin{scope}[scale=.6666667,draw=black,text=black]\node[inner sep=0pt,font=\normalsize,text=black] at (333.000,881.000) {底};\end{scope}
\begin{scope}[scale=.6666667,draw=black,text=black]\node[inner sep=0pt,font=\normalsize,text=black] at (546.000,842.000) {底};\end{scope}
\begin{scope}[scale=.6666667,draw=black,text=black]\node[inner sep=0pt,font=\normalsize,text=black] at (662.000,839.000) {底};\end{scope}
\hdsection{4}{923}{画一个以$BC$为底、$AD$为高的平行四边形。}
\begin{scope}[scale=.6666667,draw=black,text=black]\draw[hdgrid] (241,967) rectangle (673,1133);\foreach \i in {1,...,20}{\draw[hdgrid] ({241+20.571428571428573*\i},967)--({241+20.571428571428573*\i},1133);}\foreach \j in {1,...,7}{\draw[hdgrid] (241,{967+20.75*\j})--(673,{967+20.75*\j});}\end{scope}
\begin{scope}[scale=.6666667,draw=black,text=black]\draw[hdoutline] (382.000,1071.000)--(506.000,1071.000);\end{scope}
\begin{scope}[scale=.6666667,draw=black,text=black]\draw[hdoutline,dashed] (424.000,1009.000)--(424.000,1071.000);\end{scope}
\begin{scope}[scale=.6666667,draw=black,text=black]\draw[line width=.45pt] (424.000,1063.000)--(432.000,1063.000)--(432.000,1071.000);\end{scope}
\begin{scope}[scale=.6666667,draw=black,text=black]\node[inner sep=0pt,font=\normalsize,text=black] at (424.000,991.000) {$A$};\end{scope}
\begin{scope}[scale=.6666667,draw=black,text=black]\node[inner sep=0pt,font=\normalsize,text=black] at (382.000,1088.000) {$B$};\end{scope}
\begin{scope}[scale=.6666667,draw=black,text=black]\node[inner sep=0pt,font=\normalsize,text=black] at (506.000,1088.000) {$C$};\end{scope}
\begin{scope}[scale=.6666667,draw=black,text=black]\node[inner sep=0pt,font=\normalsize,text=black] at (424.000,1088.000) {$D$};\end{scope}

\begin{scope}[color=answerred,draw=answerred,text=answerred]\ans{190.00}{226.00}{106.67}{11}{①④⑤}
\ans{406.00}{226.00}{100.00}{11}{②③⑥}
\draw[answerred,line width=.8pt] (99.333,450.667) -- (177.333,450.667) -- (208.667,386.000) -- (130.667,386.000) -- cycle;
\draw[answerred,line width=.7pt,dashed] (130.667,386.000) -- (130.667,450.667);
\draw[answerred,line width=.7pt] (136.000,450.667) -- (136.000,445.333);
\draw[answerred,line width=.7pt] (136.000,445.333) -- (130.667,445.333);
\ans{134.67}{413.00}{20.00}{9}{高}
\ans{136.00}{452.67}{66.67}{9}{$\vphantom{0}5\mathrm{cm}$}
\draw[answerred,line width=.8pt] (301.333,450.000) -- (411.333,450.000) -- (443.333,385.333) -- (333.333,385.333) -- cycle;
\draw[answerred,line width=.7pt,dashed] (333.333,385.333) -- (333.333,450.000);
\draw[answerred,line width=.7pt] (338.667,450.000) -- (338.667,444.667);
\draw[answerred,line width=.7pt] (338.667,444.667) -- (333.333,444.667);
\ans{337.33}{412.33}{20.00}{9}{高}
\ans{356.00}{452.67}{66.67}{9}{$\vphantom{0}7\mathrm{cm}$}
\draw[answerred,line width=.7pt,dashed] (135.333,573.333) -- (91.366,553.041);
\draw[answerred,line width=.7pt] (89.131,557.883) -- (93.973,560.118);
\draw[answerred,line width=.7pt] (93.973,560.118) -- (96.208,555.276);
\ans{117.35}{557.85}{20.00}{9}{高}
\draw[answerred,line width=.7pt,dashed] (221.333,521.333) -- (221.333,576.667);
\draw[answerred,line width=.7pt] (226.667,576.667) -- (226.667,571.333);
\draw[answerred,line width=.7pt] (226.667,571.333) -- (221.333,571.333);
\ans{225.33}{543.67}{20.00}{9}{高}
\draw[answerred,line width=.7pt,dashed] (314.000,574.000) -- (352.000,574.000);
\draw[answerred,line width=.7pt] (352.000,568.667) -- (346.667,568.667);
\draw[answerred,line width=.7pt] (346.667,568.667) -- (346.667,574.000);
\ans{337.00}{568.67}{20.00}{9}{高}
\draw[answerred,line width=.7pt,dashed] (503.333,538.667) -- (457.659,561.073);
\draw[answerred,line width=.7pt] (455.310,556.285) -- (460.098,553.936);
\draw[answerred,line width=.7pt] (460.098,553.936) -- (462.447,558.724);
\ans{484.50}{544.54}{20.00}{9}{高}
\draw[answerred,line width=.8pt] (254.667,714.000) -- (337.333,714.000) -- (365.333,672.667) -- (282.667,672.667) -- cycle;
\draw[answerred,line width=.7pt,dashed] (282.667,672.667) -- (282.667,714.000);
\ans{286.5}{688.0}{20}{9}{高}
\ans{366.67}{656.67}{23.33}{9}{$\vphantom{0}E$}
\end{scope}
\end{scope}\pgfresetboundingbox\path[use as bounding box] (55.00,519.00) rectangle (554.00,609.00);
\end{tikzpicture}}
\end{center}\par\endgroup
\keyheading{第4题　以BC为底、AD为高的示例}
\begingroup\color{answerred}\par\noindent\begin{center}{%
\begin{tikzpicture}[x=1bp,y=-1bp,color=black,draw=black,text=black]
\begin{scope}\clip (140.00,650.00) rectangle (458.00,784.00);
\hdhead{44}{多边形的认识——信息窗3}
\hdsection{1}{145}{分一分。（填序号。）}
\begin{scope}[scale=.6666667,draw=black,text=black]\draw[hdoutline] (108.000,229.000)--(197.000,229.000)--(218.000,274.000)--(129.000,274.000)--cycle;\end{scope}
\begin{scope}[scale=.6666667,draw=black,text=black]\draw[hdoutline] (249.000,273.000)--(257.000,228.000)--(318.000,228.000)--(344.000,273.000)--cycle;\end{scope}
\begin{scope}[scale=.6666667,draw=black,text=black]\draw[hdoutline] (375.000,272.000)--(391.000,227.000)--(445.000,227.000)--(463.000,272.000)--cycle;\end{scope}
\begin{scope}[scale=.6666667,draw=black,text=black]\draw[hdoutline] (498.000,271.000)--(517.000,226.000)--(594.000,226.000)--(575.000,271.000)--cycle;\end{scope}
\begin{scope}[scale=.6666667,draw=black,text=black]\draw[hdoutline] (630.000,232.000)--(660.000,194.000)--(690.000,232.000)--(660.000,270.000)--cycle;\end{scope}
\begin{scope}[scale=.6666667,draw=black,text=black]\draw[hdoutline] (723.000,225.000)--(800.000,225.000)--(794.000,269.000)--(752.000,269.000)--cycle;\end{scope}
\begin{scope}[scale=.6666667,draw=black,text=black]\node[inner sep=0pt,font=\normalsize,text=black] at (163.000,311.000) {①};\end{scope}
\begin{scope}[scale=.6666667,draw=black,text=black]\node[inner sep=0pt,font=\normalsize,text=black] at (296.000,311.000) {②};\end{scope}
\begin{scope}[scale=.6666667,draw=black,text=black]\node[inner sep=0pt,font=\normalsize,text=black] at (419.000,311.000) {③};\end{scope}
\begin{scope}[scale=.6666667,draw=black,text=black]\node[inner sep=0pt,font=\normalsize,text=black] at (547.000,311.000) {④};\end{scope}
\begin{scope}[scale=.6666667,draw=black,text=black]\node[inner sep=0pt,font=\normalsize,text=black] at (660.000,311.000) {⑤};\end{scope}
\begin{scope}[scale=.6666667,draw=black,text=black]\node[inner sep=0pt,font=\normalsize,text=black] at (764.000,311.000) {⑥};\end{scope}
\q{144}{342}{132}{13}{平行四边形：}
\q{262}{339}{28}{13}{（}
\q{459}{339}{25}{13}{）}
\q{527}{342}{80}{13}{梯形：}
\q{583}{339}{28}{13}{（}
\q{760}{339}{25}{13}{）}
\hdsection{2}{403}{按要求画一画。}
\q{104}{447}{734}{12}{在方格纸上分别画一个底是5厘米和一个底是7厘米的平行四边形，}
\q{104}{489}{734}{12}{并各画出一条高。（每个小方格的边长为1厘米。）}
\begin{scope}[scale=.6666667,draw=black,text=black]\draw[hdgrid] (103,531) rectangle (813,700);\foreach \i in {1,...,29}{\draw[hdgrid] ({103+23.666666666666668*\i},531)--({103+23.666666666666668*\i},700);}\foreach \j in {1,...,6}{\draw[hdgrid] (103,{531+24.142857142857142*\j})--(813,{531+24.142857142857142*\j});}\end{scope}
\hdsection{3}{736}{分别画出下面平行四边形底边上的高。}
\begin{scope}[scale=.6666667,draw=black,text=black]\draw[hdoutline] (123.000,860.000)--(147.000,808.000)--(227.000,808.000)--(203.000,860.000)--cycle;\end{scope}
\begin{scope}[scale=.6666667,draw=black,text=black]\draw[hdoutline] (307.000,865.000)--(332.000,782.000)--(386.000,782.000)--(361.000,865.000)--cycle;\end{scope}
\begin{scope}[scale=.6666667,draw=black,text=black]\draw[hdoutline] (471.000,785.000)--(528.000,806.000)--(528.000,882.000)--(471.000,861.000)--cycle;\end{scope}
\begin{scope}[scale=.6666667,draw=black,text=black]\draw[hdoutline] (670.000,808.000)--(755.000,808.000)--(781.000,861.000)--(696.000,861.000)--cycle;\end{scope}
\begin{scope}[scale=.6666667,draw=black,text=black]\node[inner sep=0pt,font=\normalsize,text=black] at (118.000,830.000) {底};\end{scope}
\begin{scope}[scale=.6666667,draw=black,text=black]\node[inner sep=0pt,font=\normalsize,text=black] at (333.000,881.000) {底};\end{scope}
\begin{scope}[scale=.6666667,draw=black,text=black]\node[inner sep=0pt,font=\normalsize,text=black] at (546.000,842.000) {底};\end{scope}
\begin{scope}[scale=.6666667,draw=black,text=black]\node[inner sep=0pt,font=\normalsize,text=black] at (662.000,839.000) {底};\end{scope}
\hdsection{4}{923}{画一个以$BC$为底、$AD$为高的平行四边形。}
\begin{scope}[scale=.6666667,draw=black,text=black]\draw[hdgrid] (241,967) rectangle (673,1133);\foreach \i in {1,...,20}{\draw[hdgrid] ({241+20.571428571428573*\i},967)--({241+20.571428571428573*\i},1133);}\foreach \j in {1,...,7}{\draw[hdgrid] (241,{967+20.75*\j})--(673,{967+20.75*\j});}\end{scope}
\begin{scope}[scale=.6666667,draw=black,text=black]\draw[hdoutline] (382.000,1071.000)--(506.000,1071.000);\end{scope}
\begin{scope}[scale=.6666667,draw=black,text=black]\draw[hdoutline,dashed] (424.000,1009.000)--(424.000,1071.000);\end{scope}
\begin{scope}[scale=.6666667,draw=black,text=black]\draw[line width=.45pt] (424.000,1063.000)--(432.000,1063.000)--(432.000,1071.000);\end{scope}
\begin{scope}[scale=.6666667,draw=black,text=black]\node[inner sep=0pt,font=\normalsize,text=black] at (424.000,991.000) {$A$};\end{scope}
\begin{scope}[scale=.6666667,draw=black,text=black]\node[inner sep=0pt,font=\normalsize,text=black] at (382.000,1088.000) {$B$};\end{scope}
\begin{scope}[scale=.6666667,draw=black,text=black]\node[inner sep=0pt,font=\normalsize,text=black] at (506.000,1088.000) {$C$};\end{scope}
\begin{scope}[scale=.6666667,draw=black,text=black]\node[inner sep=0pt,font=\normalsize,text=black] at (424.000,1088.000) {$D$};\end{scope}

\begin{scope}[color=answerred,draw=answerred,text=answerred]\ans{190.00}{226.00}{106.67}{11}{①④⑤}
\ans{406.00}{226.00}{100.00}{11}{②③⑥}
\draw[answerred,line width=.8pt] (99.333,450.667) -- (177.333,450.667) -- (208.667,386.000) -- (130.667,386.000) -- cycle;
\draw[answerred,line width=.7pt,dashed] (130.667,386.000) -- (130.667,450.667);
\draw[answerred,line width=.7pt] (136.000,450.667) -- (136.000,445.333);
\draw[answerred,line width=.7pt] (136.000,445.333) -- (130.667,445.333);
\ans{134.67}{413.00}{20.00}{9}{高}
\ans{136.00}{452.67}{66.67}{9}{$\vphantom{0}5\mathrm{cm}$}
\draw[answerred,line width=.8pt] (301.333,450.000) -- (411.333,450.000) -- (443.333,385.333) -- (333.333,385.333) -- cycle;
\draw[answerred,line width=.7pt,dashed] (333.333,385.333) -- (333.333,450.000);
\draw[answerred,line width=.7pt] (338.667,450.000) -- (338.667,444.667);
\draw[answerred,line width=.7pt] (338.667,444.667) -- (333.333,444.667);
\ans{337.33}{412.33}{20.00}{9}{高}
\ans{356.00}{452.67}{66.67}{9}{$\vphantom{0}7\mathrm{cm}$}
\draw[answerred,line width=.7pt,dashed] (135.333,573.333) -- (91.366,553.041);
\draw[answerred,line width=.7pt] (89.131,557.883) -- (93.973,560.118);
\draw[answerred,line width=.7pt] (93.973,560.118) -- (96.208,555.276);
\ans{117.35}{557.85}{20.00}{9}{高}
\draw[answerred,line width=.7pt,dashed] (221.333,521.333) -- (221.333,576.667);
\draw[answerred,line width=.7pt] (226.667,576.667) -- (226.667,571.333);
\draw[answerred,line width=.7pt] (226.667,571.333) -- (221.333,571.333);
\ans{225.33}{543.67}{20.00}{9}{高}
\draw[answerred,line width=.7pt,dashed] (314.000,574.000) -- (352.000,574.000);
\draw[answerred,line width=.7pt] (352.000,568.667) -- (346.667,568.667);
\draw[answerred,line width=.7pt] (346.667,568.667) -- (346.667,574.000);
\ans{337.00}{568.67}{20.00}{9}{高}
\draw[answerred,line width=.7pt,dashed] (503.333,538.667) -- (457.659,561.073);
\draw[answerred,line width=.7pt] (455.310,556.285) -- (460.098,553.936);
\draw[answerred,line width=.7pt] (460.098,553.936) -- (462.447,558.724);
\ans{484.50}{544.54}{20.00}{9}{高}
\draw[answerred,line width=.8pt] (254.667,714.000) -- (337.333,714.000) -- (365.333,672.667) -- (282.667,672.667) -- cycle;
\draw[answerred,line width=.7pt,dashed] (282.667,672.667) -- (282.667,714.000);
\ans{286.5}{688.0}{20}{9}{高}
\ans{366.67}{656.67}{23.33}{9}{$\vphantom{0}E$}
\end{scope}
\end{scope}\pgfresetboundingbox\path[use as bounding box] (140.00,650.00) rectangle (458.00,784.00);
\end{tikzpicture}}
\end{center}\par\endgroup

\begingroup\color{answerred}\textbf{多解核对：}第2题大小和倾斜程度可不同，但底长必须分别为5厘米、7厘米。第3题平行底边之间的高可画在不同位置；第4题图形也不唯一，必须以BC为底并以AD为高。\par\endgroup
